\documentclass[12pt,a4paper]{ctexart}
\usepackage{geometry}\geometry{a4paper,margin=2.5cm,headheight=15pt}
\usepackage{fancyhdr}\pagestyle{fancy}
\fancyhead[L]{dgaiot 技术白皮书}
\fancyhead[R]{V2.0}
\fancyfoot[C]{\thepage}
\usepackage{hyperref}\hypersetup{colorlinks=true,linkcolor=blue,urlcolor=cyan}
\usepackage{listings}\lstset{basicstyle=\ttfamily\small,breaklines=true,frame=single}
\usepackage{longtable,booktabs,enumitem}

\title{{\Huge\heiti dgaiot 技术白皮书}\\{\Large 本体驱动 AI 编程 · DLAS 架构 · 确定性工业物联网}}
\author{dgaiot 项目组\\\small 2026年7月}
\date{}

\begin{document}
\maketitle
\thispagestyle{empty}
\newpage
\tableofcontents
\newpage

\section{核心设计原则}

\subsection{本体是唯一真相源}

本体（Ontology）定义了工业系统的完整语义模型——从物理设备到业务规则的映射关系。在 dgaiot 架构中，本体包括：

\begin{itemize}[leftmargin=*]
  \item 物模型（Thing Model）：125 个 Modbus 测点定义，含地址、类型、阈值
  \item 产品模板（ProductTemplet）：设备类型、属性、事件、服务
  \item 四层拓扑（Site/Gateway/Device/Point）：物理世界到数字孪生的映射
  \item 通道配置（Channel）：协议适配器与 TDengine 数据库映射
  \item ACL 规则：设备授权（ProductSecret/DeviceSecret）与用户角色权限
\end{itemize}

\subsection{生产环境运行确定性代码}

\begin{quote}
开发时 AI 编程（读本体 $\rightarrow$ 写代码）。\\
运行时确定性执行（gen\_statem $\cdot$ MQTT $\cdot$ TDengine）。\\
没有 AI 推理，没有概率，没有黑盒。\\
每一步可追溯、可回放、可审计。
\end{quote}

\subsection{修正本体，而非修正代码}

本体是严肃的——企业标准、合规要求、物理规律。代码是派生的。

\begin{itemize}[leftmargin=*]
  \item 日志显示温度 82\textdegree C 未告警 $\rightarrow$ 检查 gen\_statem 是否收到消息（运行时问题）
  \item 再检查 condition 是否正确编译（AI 生成问题）
  \item 最后才怀疑本体（极少）
  \item 本体 = 宪法，不常修。代码 = 法规，本体派生，可迭代。
\end{itemize}

\section{DLAS 架构}

\subsection{五层架构}

\begin{lstlisting}[language=]
EDGE:    iotStudio (Python+Vue) — DeviceAccess(9) → StreamEngine(15)
SECURITY: auth · role · ACL/CLP · Hooks · JWT
ACTION:  Shadow(gen_statem) · Bridge · MQTT · Rule Engine
LOGIC:   Ontology Engine · Model Registry · Reasoner · Policy Optimizer
DATA:    Parse/PG(23类) · TDengine(_{ChannelId}._{ProductId}) · EMQX(:1883)
\end{lstlisting}

\subsection{影子设备：物理世界 1:1 映射}

每个物理设备对应一个 Erlang OTP gen\_statem 进程：

\begin{lstlisting}[language=erlang]
init() -> {ok, authenticate, Device}.
authenticate(cast, heartbeat) -> {next_state, online, Device}.
online(cast, {data, Props}, Device) ->
    evaluate(Rules, Props),  %% 物模型编译的规则
    bridge -> Parse + TDengine,
    case alarm_triggered of
        true  -> {next_state, alarm, Device};
        false -> {keep_state, Device}
    end.
alarm(cast, heartbeat, Device) ->
    {next_state, online, Device}.  %% 恢复
\end{lstlisting}

211 台设备的数字孪生 = 211 个 Erlang 进程，总内存消耗 $< 1$ MB。

\subsection{TDengine 存储结构}

\begin{lstlisting}[language=SQL]
-- 宏: ?Database(Name) = "_" ++ Name
-- 宏: ?Table(Name)    = "_" ++ Name

CREATE DATABASE IF NOT EXISTS _85ef6b7459 KEEP 365;
CREATE STABLE IF NOT EXISTS _85ef6b7459._2de1b3e1b8
  (ts TIMESTAMP, value FLOAT, quality INT)
  TAGS (devaddr NCHAR(50));
CREATE TABLE ... USING _2de1b3e1b8 TAGS('DEV-001');

INSERT INTO _85ef6b7459.sub_2de1b3e1b8_DEV-001
  VALUES (NOW, 2.35, 192);
\end{lstlisting}

\subsection{DLAS 数据流}

\begin{lstlisting}[language=]
Modbus(:502) -> Edge Collector -> MQTT
  -> EMQX -> Shadow gen_statem -> evaluate(Rules)
  -> Parse(update status) + TDengine(INSERT telemetry)
\end{lstlisting}

\section{本体驱动 AI 编程}

\subsection{闭环流程}

\begin{enumerate}[leftmargin=*]
  \item 企业需求：合规、SLA、安全策略、业务规则
  \item 本体建模：thing\_model、ProductTemplet、Channel、ACL
  \item AI 生成代码：.erl（gen\_statem + rules）、.py（bridge）、SQL（schema）
  \item 测试验证：单元测试（state transitions）、集成测试（MQTT chain）、日志验证（TDengine INSERT）
  \item 确定性生产：gen\_statem evaluate、dgiot\_tdengine:INSERT、MQTT publish
  \item 日志反馈：错误日志 $\rightarrow$ 修正 AI 生成逻辑；性能日志 $\rightarrow$ 优化参数
\end{enumerate}

\subsection{本体 $\rightarrow$ 代码 示例}

\begin{lstlisting}[language=json, caption=thing\_model.json 片段]
{
  "name": "油压",
  "dataForm": {"address":"40300", "protocol":"modbus",
               "originaltype":"float32_AB"},
  "identifier": "oil_pressure"
}
\end{lstlisting}

AI 自动生成：

\begin{lstlisting}[language=erlang]
%% Modbus: readHoldingRegisters(40300, 2, float32_AB)
%% MQTT: dgiot/oil_field_01/gw_131/rtu_001/oil_pressure/data
%% TDengine: INSERT INTO _85ef6b7459._2de1b3e1b8 VALUES (NOW, v, 192)
%% gen_statem:
online(cast, {data, #{oil_pressure := V}}, Device)
  when V > 3.0 ->
    {next_state, alarm, Device, [{state_timeout, 30000, heartbeat_missed}]};
\end{lstlisting}

\section{授权模型}

\subsection{三层 MQTT 授权}

\begin{table}[H]\centering
\caption{设备/用户/超级用户三层授权}
\begin{tabular}{p{3cm} p{3cm} p{4cm} p{3cm}}
\toprule
\textbf{层级} & \textbf{ClientID} & \textbf{认证方式} & \textbf{权限范围} \\
\midrule
设备授权 & \{PID\}\_\{DevAddr\} & ProductSecret/DeviceSecret & 上报遥测 \\
用户授权 & \{Token\}\{Type\} & Session Token + Role ACL & 订阅/下发 \\
超级用户 & dgiot (127.0.0.1) & 本地免检 & 内部通信 \\
\bottomrule
\end{tabular}
\end{table}

\section{验证结果}

\subsection{全链路测试}

\begin{table}[H]\centering
\caption{端到端验证矩阵}
\begin{tabular}{p{4cm} p{3cm} p{5cm}}
\toprule
\textbf{测试项} & \textbf{方法} & \textbf{结果} \\
\midrule
MQTT pub/sub & Kylin-DMZ 实机 & 3发3收，DLAS topic 格式 \\
gen\_statem Shadow & Erlang 编译运行 & init$\rightarrow$online$\rightarrow$alarm$\rightarrow$online \\
TDengine 存储 & 源码对齐建表 & INSERT 3行，SELECT 3行 \\
dlink 设备授权 & ClientID=\{PID\}\_\{DevAddr\} & 9点写入，全部 OK \\
\bottomrule
\end{tabular}
\end{table}

\section{结论}

dgaiot 提出了一种\textbf{本体驱动 AI 编程}的工业物联网架构。本体是唯一真相源，AI 在开发阶段将本体编译为确定性代码，生产环境运行的是 gen\_statem + MQTT + TDengine 的确定性管道，不包含任何 AI 推理或概率决策。DLAS 五层架构已通过 Kylin-DMZ 实测验证。

\end{document}
